\documentclass[11pt]{article}
\usepackage[a4paper,margin=1in]{geometry}
\usepackage{setspace}
\usepackage{booktabs}
\usepackage{amsmath,amssymb}
\usepackage{graphicx}
\usepackage{hyperref}
\usepackage{float}
\onehalfspacing

\title{Geospatial Climate Data and Economic Outcomes: \\
A State-Month Panel for India using ERA5 Precipitation}
\author{Your Name}
\date{\today}

\begin{document}
\maketitle

\section{Research Question and Motivation}
This project studies whether short-run rainfall shocks, measured from ERA5 reanalysis data,
affect a state-level economic outcome in India. The assignment objective is to combine:
(i) geospatial raster processing, (ii) zonal aggregation to administrative units, and
(iii) panel econometrics with fixed effects.

The economic interpretation is straightforward: rainfall fluctuations can influence agricultural
production, food supply chains, and inflation dynamics. At monthly frequency, precipitation
anomalies can generate meaningful transitory variation relevant for macroeconomic outcomes.

\section{Data}
\subsection{Climate data}
The climate variable is ERA5 \texttt{total\_precipitation} at daily frequency over India.
Administrative units are India states (GADM level-1 shapefiles). The analysis period is
2015--2024, with 2015--2023 as climatology baseline.

\subsection{Geospatial processing}
The pipeline performs the following steps:
\begin{enumerate}
    \item Read ERA5 netCDF as a raster stack.
    \item Read state boundaries and harmonize coordinate reference systems (CRS).
    \item Compute zonal means of daily precipitation in each state.
    \item Convert precipitation from meters to millimeters.
\end{enumerate}

\section{Climate Shock Construction}
Let $p_{itd}$ denote daily precipitation (mm) in state $i$ on day $d$ in month $t$.
Monthly mean precipitation is:
\begin{equation}
P_{it} = \frac{1}{D_t}\sum_{d=1}^{D_t} p_{itd},
\end{equation}
where $D_t$ is the number of days in month $t$.

Two shock measures are constructed:
\begin{align}
\text{Anomaly}_{it} &= P_{it} - \bar{P}_{i,m}, \\
\text{zShock}_{it} &= \frac{P_{it} - \bar{P}_{i,m}}{\sigma_{i,m}},
\end{align}
where $\bar{P}_{i,m}$ and $\sigma_{i,m}$ are the long-run mean and standard deviation for
state $i$ in calendar month $m$.

As a robustness measure, heavy rainfall days are also computed:
\begin{equation}
HR_{it} = \sum_{d=1}^{D_t} \mathbf{1}(p_{itd} > 50\,\text{mm}).
\end{equation}

\section{Economic Outcome and Panel Merge}
The code expects a monthly state-level economic dataset with columns:
\texttt{unit\_id, ym, y\_outcome}. This can represent food inflation, agricultural output,
or another credible public indicator at comparable frequency and geography.

Climate and economic datasets are merged on state identifier and month.

\section{Empirical Strategy}
The baseline fixed-effects model is:
\begin{equation}
Y_{it} = \beta\,\text{Shock}_{it} + \alpha_i + \gamma_t + \varepsilon_{it},
\end{equation}
where $\alpha_i$ are state fixed effects and $\gamma_t$ are month-time fixed effects.
Standard errors are clustered at the state level.

To study dynamics, local-projection style regressions are estimated:
\begin{equation}
Y_{i,t+h} = \beta_h\,\text{Shock}_{it} + \alpha_i + \gamma_t + \varepsilon_{i,t+h},
\quad h=0,\dots,6.
\end{equation}

\section{Implementation Notes and Reproducibility}
All code is provided in \texttt{code/climate\_econ\_assignment.R}. Outputs are written to
\texttt{outputs/tables/} and \texttt{outputs/figures/}. A small example cleaned file is included in
\texttt{data/clean/economic\_outcome\_monthly.csv} to illustrate expected schema.

\section{Interpretation Guide for Pre-Doc / Early PhD Level}
When reporting results, emphasize:
\begin{itemize}
    \item \textbf{Identification margin:} variation is within-state over time after removing common
    aggregate shocks through time fixed effects.
    \item \textbf{Economic magnitude:} convert coefficients into interpretable units (e.g., effect of
    a one-standard-deviation rainfall shock).
    \item \textbf{Robustness:} compare anomaly and z-score definitions, and optionally include
    heavy-rainfall-day controls.
    \item \textbf{Limitations:} discuss measurement error in gridded climate data and possible
    omitted channels that co-move with weather.
\end{itemize}

\section{Conclusion}
This assignment delivers a reproducible framework connecting geospatial climate data and
macroeconomic outcomes in panel form. The workflow is directly extensible to district-level
analysis, richer outcomes, and advanced designs (e.g., dynamic treatment effects, state-specific
trends, or nonlinear precipitation impacts).

\end{document}
